\documentclass[12pt,a4paper]{report}

\usepackage[utf8]{inputenc}
\usepackage[T1]{fontenc}
\usepackage[spanish]{babel}
\usepackage{geometry}
\usepackage{graphicx}
\usepackage{float}
\usepackage{array}
\usepackage{amsmath}
\usepackage{hyperref}

\geometry{margin=2.5cm}

\begin{document}
≈
\section{Conclusion}

La version final del procesador RISC-V pipeline amplia el conjunto de instrucciones respecto a \texttt{pipev1}, incorporando nuevas operaciones aritmeticas, logicas, desplazamientos, branches, \texttt{jalr} y \texttt{lui}. Ademas, se implemento una Hazard Unit capaz de resolver dependencias mediante forwarding, stalling y flushing.

Las pruebas con programas sin dependencias, con forwarding, con load-use hazard y con flush demuestran que el datapath y la unidad de control funcionan correctamente. La comparacion con la version sin Hazard Unit evidencia que esta unidad es necesaria para evitar resultados incorrectos en presencia de dependencias de datos y hazards de control.

\end{document}
